% FIG ch11-01 three levels (Chinese)
\usetikzlibrary{arrows.meta,positioning,calc}
\begin{tikzpicture}[x=\linewidth/100,y=\linewidth/100]
\begin{scope}[
  font=\sffamily\small,
  level/.style={
    draw=BookRule, line width=0.5pt, rounded corners=2pt,
    minimum width=52mm, minimum height=11mm, align=left, inner sep=4pt,
    text width=50mm, anchor=west,
  },
  qbox/.style={
    draw=BookRule, fill=BookPale!40, rounded corners=2pt,
    align=center, inner sep=3pt, text width=22mm, font=\scriptsize,
  },
]
  \node[
    draw=BookTeal, line width=0.7pt, fill=white, rounded corners=3pt,
    align=center, inner sep=5pt, text width=24mm, font=\small,
  ] (bubble) at (8,42) {再好看一点};
  \draw[BookInk!55, dashed, -{Stealth[length=2mm]}] (bubble.east) -- (22,42);

  \node[qbox] (q1) at (32,58) {视线\\停在哪？};
  \node[qbox, below=3mm of q1] (q2) {读者会\\误解哪里？};
  \node[qbox, below=3mm of q2] (q3) {哪些本应\\一样却不一样？};

  \foreach \q in {q1,q2,q3} {
    \draw[BookInk!45, -{Stealth[length=1.6mm]}] (\q.east) -- ++(6,0);
  }

  \node[level, fill=BookInk, text=white] (L1) at (54,66) {%
    \textbf{第 1 层　科学错误}\\[-1pt]
    {\scriptsize 例：画成了悬索桥；数据箭头方向反了}\\
    {\scriptsize\textcolor{BookPale!90}{处理：回到任务书，必要时重新生成}}};
  \node[level, fill=BookInk!60, text=BookInk, below=2.5mm of L1] (L2) {%
    \textbf{第 2 层　阅读障碍}\\[-1pt]
    {\scriptsize 例：无线链路像拉索；160\,mm 下图标太小}\\
    {\scriptsize 处理：一次只改一类问题}};
  \node[level, fill=BookInk!30, text=BookInk, below=2.5mm of L2] (L3) {%
    \textbf{第 3 层　视觉一致性}\\[-1pt]
    {\scriptsize 例：同类传感器颜色不同；线宽不一}\\
    {\scriptsize 处理：优先在排版层统一}};

  \coordinate (ordA) at ($(L1.west)+(-3,0)$);
  \coordinate (ordB) at ($(L3.west)+(-3,0)$);
  \draw[BookInk!70, line width=0.8pt, -{Stealth[length=2.4mm]}] (ordA) -- (ordB);
  \node[font=\scriptsize, align=center, text width=46mm, anchor=north]
    at ($(L3.south)+(0,-2.5mm)$) {按此顺序处理；每改一次，整图复核三层};
\end{scope}
\end{tikzpicture}
